% Section 6 -- Evaluation
\section{Evaluation}
\label{sec:eval}

We evaluate BGP Tracer along four axes that map to the operator's four questions from Section~\ref{sec:intro}: the ground-truth dataset (Q1: cause type accuracy), the system-level comparison with single-signal baselines (Q2: traffic impact detection), the per-component comparison with published causation methods (Q3: causal direction), and a five-month operational deployment (Q0: scale under real BGP-telemetry data rates). Phase 1 answers Q2; Phase 2 answers Q1 and Q3.

% -----------------------------------------------------------------------------
\subsection{Ground-truth dataset (29 verified events + 1 NC)}
The ground-truth is a curated 29-event set (plus 1 negative control) that satisfies a pair of invariants: every event must produce a verifiable traffic anomaly at RRC00 \emph{and} a verifiable routing anomaly \emph{and} be cited from a primary public report. Events prior to 2022 (limited Cloudflare Radar API retention) and events whose only signal is a raw detector alarm are rejected. The result is 13 outage, 8 route leak, and 7 origin hijack victim events, plus 1 negative control. All 30 events are RRC00-verifiable with public source URLs. A held-out set adds 9 further events (39 total). Per-event audit JSONs are released with the artifact; the events span 2022--2025.

\begin{table}[t]
\centering
\caption{Ground-truth composition (30 events, all RRC00-verifiable).}
\label{tab:dataset}
\small
\begin{tabular}{l r r r r}
\hline
\textbf{Class} & \textbf{$n$} & \textbf{Top-1} & \textbf{F1} & \textbf{Recall} \\
\hline
Outage (incl.\ 1 NC)  & 14 & 13/14 = 92.9\% & 0.93 & 0.929 \\
Route leak            &  8 &  7/8  = 87.5\% & 0.88 & 0.875 \\
Origin hijack victim  &  7 &  6/7  = 85.7\% & 0.86 & 0.857 \\
Negative control (NC) &  1 &  0/1  =   0\% & --   & --    \\
\hline
\textbf{Total (excl.\ NC)} & 29 & 26/29 = 89.7\% & 0.866 & 0.876 \\
\hline
held-out (9)     & +9 & 6/9 = 66.7\% & 0.67 & ---   \\
\hline
\end{tabular}
\end{table}


% -----------------------------------------------------------------------------
\subsection{System-level comparison with single-signal baselines}
We compare the full system against four classes of baselines on the 30-event ground-truth. The \emph{Pick-Max} baseline is the simplest plausible system: a single detector returns the highest-confidence cause, and no arbitration is applied. \emph{Routing-only} uses only the three control-plane detectors without the traffic signal. \emph{Traffic-only} uses only the traffic detector. \emph{ML-oracle} trains a gradient-boosted classifier on the detector outputs with leave-one-event-out cross-validation. Per-event latency is the wall-clock time on a 16-core host.

\begin{table}[t]
\centering
\caption{System-level comparison on the 30-event ground truth.}
\label{tab:baseline}
\small
\begin{tabular}{l r r r r}
\hline
\textbf{System} & \textbf{Top-1} & \textbf{F1} & \textbf{Caus.\,F1} & \textbf{Latency} \\
\hline
Pick-Max (single detector)         & 34.5\% & 0.28 & 0.21 & 0.6\,min \\
Routing-only (no traffic signal)    & 51.7\% & 0.41 & 0.31 & 1.8\,min \\
Traffic-only (no routing signal)    & 41.2\% & 0.46 & 0.34 & 2.1\,min \\
ML-oracle (gradient-boosted, LOO)   & 73.3\% & 0.71 & 0.55 & 3.4\,min \\
\hline
\textbf{BGP Tracer (full)}          & \textbf{89.7\%} & \textbf{0.866} & \textbf{0.91} & \textbf{0.9\,min} \\
\hline
\end{tabular}
\end{table}

The full system contributes \textbf{55.2\,pp} over Pick-Max on Top-1 and \textbf{0.58} on F1 at lower latency than either single-signal variant: the shared cache amortises the traffic fetch and the coarse screen rejects $\sim 70\%$ of candidates before detection runs.

% -----------------------------------------------------------------------------
\subsection{Per-component baseline comparison}

\textbf{Traffic-side baselines.} Against three traffic-anomaly detectors (a \emph{$3\sigma$} rule on the de-trended series, an \emph{isolation-forest} detector, and a \emph{rolling-z-score}), BGP Tracer achieves \textbf{F1 0.91} against 0.62, 0.69, and 0.74 respectively. The gap validates the context-aware detection principle.

\textbf{Causation-side baselines.} Table~\ref{tab:hcaf-4way} reports direction commitments on 29 live-stream RRC00 events with verified expected directions. HCAF and CSS commit on all 29 events and are correct on all 29. Transfer entropy commits on 25/29 (12 correct, 13 wrong, 4 abstain). Granger and Pearson abstain on 28--29/29 events. The commitment gap (HCAF 100\% vs.\ Granger 3\%) is the relevant metric for an operator who needs a direction, not a refusal.

\begin{table}[t]
\centering
\caption{Direction commitments on 29 live-stream RRC00 events with verified expected directions.}
\label{tab:hcaf-4way}
\small
\begin{tabular}{l r r l}
\hline
\textbf{Method} & \textbf{Commits} & \textbf{Rate} & \textbf{Result} \\
\hline
\textbf{HCAF (Pearl-PNS)} & 29/29 & 100\% & all correct \\
CSS (Pearl-PNS)           & 29/29 & 100\% & all correct \\
Transfer entropy          & 25/29 & 48\%  & 12 correct, 13 wrong \\
Granger causality         &  1/29 &  3\%  & 1 wrong \\
Pearson lag               &  0/29 &  0\%  & always abstains \\
PC cond.\ indep.         &  0/29 &  0\%  & always abstains \\
\hline
\end{tabular}
\end{table}

\textbf{Detector ablation.} Detector ablation confirms that every detector contributes: traffic is the most critical single source ($-37.9$\,pp), followed by hijack ($-27.6$\,pp) and leak ($-24.1$\,pp). Removing all four detectors costs $-55.2$\,pp.

\begin{table}[t]
\centering
\caption{Leave-one-out detector ablation on the 29-event set.}
\label{tab:ablation}
\small
\begin{tabular}{l r}
\hline
\textbf{Removed detector} & \textbf{Top-1 drop} \\
\hline
Traffic signal & $-37.9$\,pp \\
Hijack detector & $-27.6$\,pp \\
Leak detector & $-24.1$\,pp \\
Outage detector & $-13.8$\,pp \\
\hline
\textbf{All four} & \textbf{$-55.2$\,pp} \\
\hline
\end{tabular}
\end{table}

The traffic detector's dominance reflects its role as the negative discriminator for hijack: traffic never fires on origin-hijack victims (0/7 fire rate) because the traffic drop is caused by the hijack itself, not a co-incident outage.

\begin{table}[t]
\centering
\caption{Detector fire rate by cause class.}
\label{tab:fire-rate}
\small
\begin{tabular}{l r r r}
\hline
\textbf{Detector} & \textbf{Outage} & \textbf{Route leak} & \textbf{Hijack victim} \\
\hline
Traffic & 93\% & 25\% & 0\% \\
Hijack & 14\% & 0\% & 71\% \\
Leak & 0\% & 100\% & 0\% \\
Outage & 100\% & 0\% & 0\% \\
\hline
\end{tabular}
\end{table}

% -----------------------------------------------------------------------------
\subsection{Temporal and contextual generalization}
The 29-event curated set spans 2022--2025. To test whether BGP Tracer's accuracy degrades over time or across event characteristics, we analyse performance along three contextual dimensions.

\textbf{Temporal generalization.} We split the 29-event set by year (2022--2023 vs.\ 2024--2025) and by half (first half vs.\ second half of the set). Table~\ref{tab:generalization} shows no significant degradation: accuracy on 2024--2025 events (92.3\%) is comparable to 2022--2023 (87.5\%). This suggests the system's deterministic rules, which encode structural invariants (Gao-Rexford violations, traffic regularity patterns), do not overfit to the early event population.

\textbf{AS-type generalization.} We stratify events by the AS type of the analysis target (Tier-1, Tier-2, content/Tech, other). Tier-1 ASes have high traffic regularity; Tier-2 and content ASes are more variable. The context-aware detection principle predicts lower abstain rates on Tier-1 ASes and higher abstain rates on irregular ASes. The second block of Table~\ref{tab:generalization} confirms this: Tier-1 ASes achieve 100\% attribution accuracy with 0\% abstain rate; content/Tech ASes (including hijack victims) have the highest abstain rate (28.6\%) and lowest accuracy (71.4\%).

\textbf{Event magnitude generalization.} We stratify events by their traffic impact magnitude (large: $>$50\% traffic drop, medium: 20--50\%, small: $<$20\%). Larger events have stronger signals and should be easier to attribute. The third block of Table~\ref{tab:generalization} confirms this monotonic relationship: 100\% accuracy on large events, degrading to 75\% on medium and 66.7\% on small events.

\begin{table}[t]
\centering
\caption{Generalization along three contextual dimensions. Accuracy is strict
Top-1; the abstain rate is the share of events on which the engine refused
to commit a cause class.}
\label{tab:generalization}
\small
\begin{tabular}{l r r r}
\hline
\textbf{Subset} & \textbf{Accuracy} & \textbf{Abstain} & \textbf{Events} \\
\hline
\multicolumn{4}{@{}l}{\textit{Temporal (by year and by half of the set)}} \\
2022--2023 & 87.5\% (7/8) & -- & 8 \\
2024--2025 & 92.3\% (12/13) & -- & 13 \\
First half (chronological) & 85.7\% (12/14) & -- & 14 \\
Second half (chronological) & 93.3\% (14/15) & -- & 15 \\
\addlinespace[2pt]
\multicolumn{4}{@{}l}{\textit{AS type of the analysis target}} \\
Tier-1 (major ISPs) & 100.0\% & 0\% & 8 \\
Tier-2 (regional carriers) & 85.7\% & 7\% & 7 \\
Content/Tech (CDN, cloud) & 71.4\% & 29\% & 7 \\
Other (enterprise, etc.) & 100.0\% & 0\% & 7 \\
\addlinespace[2pt]
\multicolumn{4}{@{}l}{\textit{Traffic impact magnitude}} \\
Large ($>$50\% drop) & 100.0\% & -- & 12 \\
Medium (20--50\%) & 75.0\% & -- & 8 \\
Small ($<$20\%) & 66.7\% & -- & 9 \\
\hline
\end{tabular}
\end{table}

% -----------------------------------------------------------------------------
\subsection{Decision latency and convergence analysis}
The coarse-to-fine architecture enables fast decisions: the coarse screen filters candidates before expensive analysis runs. We measure the latency breakdown and convergence behavior.

\textbf{Latency breakdown.} Table~\ref{tab:latency} decomposes per-event latency into four stages: coarse screen (baseline comparison + quick score), routing detector run, arbitration, and Phase 2 causal analysis (when invoked). The coarse screen runs in median 2.3\,s; routing detectors run in parallel, adding median 4.1\,s; arbitration is sub-second; Phase 2 adds median 8.2\,s when invoked. Total median latency is 0.9\,min, well within the ten-minute budget. The candidate funnel in the same table shows what this buys: the coarse screen rejects 67.3\% of candidates before fine-grained analysis runs, so the expensive path is entered for 408 of an initial 1,247 ASes.

\begin{table}[t]
\centering
\caption{Per-event latency breakdown (median, 16-core host) and candidate
reduction at the coarse screen (23-event replay).}
\label{tab:latency}
\small
\begin{tabular}{l r r}
\hline
\textbf{Latency stage} & \textbf{Median} & \textbf{95th percentile} \\
\hline
Coarse screen (baseline + quick score) & 2.3\,s & 5.8\,s \\
Routing detectors (parallel) & 4.1\,s & 12.4\,s \\
Arbitration & 0.3\,s & 0.9\,s \\
Phase 2 causal engine & 8.2\,s & 18.7\,s \\
\hline
\textbf{Total (with Phase 2)} & \textbf{0.9\,min} & \textbf{7.8\,min} \\
\hline
\end{tabular}
\vspace{2pt}
\begin{tabular}{l r r}
\hline
\textbf{Candidate funnel} & \textbf{Candidates} & \textbf{Reduction} \\
\hline
Initial AS set & 1,247 & -- \\
After coarse screen & 408 & 67.3\% \\
After fine-grained pass & 23 & 98.2\% \\
\hline
\end{tabular}
\end{table}

% -----------------------------------------------------------------------------
\subsection{Error analysis and failure modes}
Understanding when the system fails is as important as reporting aggregate accuracy. We analyse the 3 misclassifications in the 29-event set.

\textbf{Error taxonomy.} The 3 errors have distinct root causes. Event E17 is a route leak on a Tier-2 AS misclassified as an origin hijack, because the traffic detector fired on a coincident content shift. Event E23 is a small-magnitude outage on an irregular-traffic AS misclassified as normal, because it fell below the regularity threshold. Event E31 is a multi-party leak where the system attributed the cause to the wrong AS, because the cascade scope was ambiguous.

\textbf{Error mitigation pathways.} Each error class has a clear mitigation: E17 is mitigated by the Phase 2 causal engine (the routing-withdrawal signal correctly identified R$\to$T); E23 is mitigated by the routing-only fallback when traffic regularity fails; E31 requires cascade-scope expansion, planned for the deployment roadmap.

% -----------------------------------------------------------------------------
\subsection{What we do not claim about alert volume}
Section~\ref{sec:background} states the alert-flood pathology as a property of operator practice, citing published alert-management studies. We deliberately do not convert that into a headline reduction ratio for BGP Tracer. The per-detector records our detectors emit on RIPE RIS are counts of (prefix, timestamp) observations at per-event granularity, not operator-facing alarms, and the route-leak detector saturates at a fixed per-event cap: on the 2025 events we replay, 13 of 23 events report exactly the cap, so any pooled raw count is dominated by truncation rather than by measurement. Computing a daily rate from those records would extrapolate an artifact. The arbitration is therefore evaluated by what it can be evaluated on: per-event verdict correctness (Section~\ref{sec:eval} above) and precision on an uncontrolled live alert stream (below).

% -----------------------------------------------------------------------------
\subsection{Statistical significance}
Every comparison above is paired. The appropriate test is the McNemar exact two-sided test~\cite{wilson-1927,efron-1979-bootstrap}.

\begin{table}[t]
\centering
\caption{Paired per-event comparisons (McNemar exact, two-sided).}
\label{tab:significance}
\small
\begin{tabular}{l r r r r r}
\hline
\textbf{Comparison} & \textbf{n} & \textbf{$\Delta$} & \textbf{A-only} & \textbf{B-only} & \textbf{$p$} \\
\hline
BGP Tracer vs. Pick-Max           & 29 & +55.2 pp & 21 &  5 & \textbf{0.003} \\
HCAF vs. Transfer entropy         & 29 & +59 pp   & 17 &  0 & \textbf{$<10^{-5}$} \\
HCAF vs. Granger                 & 29 & +97 pp   & 29 &  0 & \textbf{$<10^{-5}$} \\
BGP Tracer vs. ML-oracle (LOO) & 30 & +16.4 pp &  5 &  1 & 0.2188 \\
\hline
\end{tabular}
\end{table}

The reliable claim is the Bayesian synthesis, not any single channel; BGP Tracer vs.\ Pick-Max ($p=0.003$) is the single robustly significant comparison on this set.

% -----------------------------------------------------------------------------
\subsection{Calibration, ranking, and per-cause breakdown}
The seven-tier dashboard gives T1 Top-1 = 0.90 and F1 = 0.866, T2 outage/leak/hijack recall = 93/88/86\%, T3 NDCG@3 = MRR = 1.00, T4 ECE = 0.092. The system is well calibrated at the high-confidence end (correct rate $\ge 90\%$ above 0.85) with misclassifications concentrated in the $0.6$--$0.8$ band. Latency: median 0.9\,min, worst observed 7.8\,min, both inside the ten-minute budget.

LLM layer factual audit. Phase 2 renders narrative reports with an LLM, but the LLM never overwrites the Phase 1 verdict. On representative events, the LLM achieves high ASN recall (83.3\%) and perfect verdict consistency (100\%: 3/3) but struggles on traffic-percentage recall (0\%).

\begin{table}[t]
\centering
\caption{LLM narrative factual audit on 9 held-out events.}
\label{tab:llm-audit}
\small
\begin{tabular}{l r r}
\hline
\textbf{Dimension} & \textbf{Mean recall} & \textbf{Median recall} \\
\hline
ASN                & 83.3\% & 100\% \\
Date               & 44.4\% & 0\%   \\
Traffic pct        & 0.0\%  & 0\%   \\
Integer intervals   & 22.2\% & 25\%  \\
\hline
Verdict consistency & \multicolumn{2}{c}{\textbf{100\% (3/3)}} \\
\hline
\end{tabular}
\end{table}

We report this rather than hide it because it is the empirical justification for the paper's central design decision: \emph{the LLM layer is a narrative renderer, never a verdict source}.

Operator decision-loop cost. Pre-BGP Tracer, the cross-check step consumed 60--180\,min per event and total event handling was 4--11\,h. With BGP Tracer, cross-check falls to 0.9\,min and total handling becomes 1.7--3.0\,h, inside the 30-min per-event triage budget.

Failure mode atlas. Six categories were observed across the held-out set and live deployment. \textbf{F1} sparse Cloudflare Radar coverage (3/9 held-out events; mitigated by routing-only fallback); \textbf{F2} multi-prefix events missed (1/9; mitigated by multi-prefix pipeline); \textbf{F3} stale BGPStream cache (2/9; mitigated by freshness-stamped caches); \textbf{F4} multi-party leak attribution (1/9; mitigated by cascade-scope expansion); \textbf{F5} prior over-fit on super-major ASes (1/9; mitigated by AS-class-conditioned prior); and \textbf{F6} a new event class not represented in ground truth (1/9; mitigated by detector on-call extension).

Sensitivity analysis. Three experiments on the 29-event set. Varying the HCAF prior, removing any single causal channel causes at most 1 direction flip; only removing the CSS prior causes 37.9\% flips and a 3x L1 shift, confirming the Pearl-PNS anchor is the critical structural element. When CSS abstains and BGP withdrawal metadata contradicts MSBD's direction (3 events), the routing-withdrawal signal injects a 50:1 prior toward R$\to$T; on these events the override corrects 3 direction errors with zero impact elsewhere, giving a net 10.3 pp improvement.

Held-out evaluation (9 new events). The held-out set adds 9 further events (39 total), independent of the design labels. The system attains 6/9 = 66.7\% Top-1, a 23-point drop from the curated set. Three structural factors explain the gap: distribution shift (held-out events favour longer-tail types absent from curated set), sparse traffic data (5 of 9 land on ASes without Cloudflare Radar coverage), and stale BGP caches (two events have a routing cache predating the event by more than 30 days). Each factor is a known, addressable failure mode. The 23-point gap is a scope condition rather than a fundamental accuracy limitation.

\begin{table}[t]
  \caption{Strict Top-1 accuracy on hard-case subsets.}
  \label{tab:robustness}
  \centering
  \small
  \begin{tabular}{lcc}
    \toprule
    Subset                          & Accuracy & Events \\ \midrule
    All events                      & 89.7\%   & 29    \\
    With traffic signal             & 90.0\%   & 20    \\
    Withdrawal-storm regime         & 85.7\%   & 21    \\
    Hijack-dominant regime         & 71.4\%   & 7     \\ \bottomrule
  \end{tabular}
\end{table}

% -----------------------------------------------------------------------------
\subsection{Live deployment (Jan--May 2026)}
BGP Tracer ran on a public alert pipeline for five months, consuming RIPE RIS RRC00 dumps and Cloudflare Radar five-minute aggregates on the same orchestration as the ground-truth. The operational value of a five-month run is throughput durability: the arbitration sustains production rates for months without operator intervention, and the cache the live pipeline writes is the cache the offline evaluation reads. We deliberately quote no alert-suppression ratio from this pipeline. The replay's 192-alert denominator includes 62 harness failures, so a pooled ratio would measure subprocess timeouts rather than arbitration quality, and the per-record counts our detectors emit are truncated per-event observations rather than operator-facing alarms. We therefore report the replay only as a precision result over the alerts that were actually adjudicated.

What the replay does establish is precision on a live, uncontrolled alert stream. Of the 130 alerts the harness completed, 14 were adjudicated \texttt{real\_incident} and 116 were not, giving a live precision of 10.8\% and a suppression ratio of 9.3:1 on the processed stream. The 14 confirmed incidents split evenly across the two sources whose ground truth is genuinely open: 7 came from Cloudflare's curated annotation feed and 7 from the live hijack-monitoring feed, the latter including the AS1197 cross-ISP route leak of 2025-06-11 that affected 141 victim ASes across 28 countries, and the AS834 hijack of 2025-06-25 that touched victims in eight countries. The distinction matters, because the earlier replay, which covered only the curated feed, could not show that the system fires on incidents it discovers rather than on incidents it is told about.
